% The article for the keyword-and-ledger pull request: kept statistics, one statement, and an
% index of the attestation ledger by output hash, measured on the corpus instance.
\documentclass[11pt]{article}
\usepackage[T1]{fontenc}
\usepackage{lmodern}
\usepackage{microtype}
\usepackage[margin=1in]{geometry}
\usepackage{amsmath,amssymb,amsthm,mathtools}
\usepackage{algorithm}
\usepackage{algpseudocode}
\usepackage{booktabs}
\usepackage{array}
\usepackage[numbers,sort&compress]{natbib}
\usepackage[hidelinks]{hyperref}
\usepackage[capitalise,noabbrev]{cleveref}

\setlength{\emergencystretch}{3em}
\renewcommand{\topfraction}{0.9}
\renewcommand{\textfraction}{0.05}
\renewcommand{\floatpagefraction}{0.75}

\theoremstyle{plain}
\newtheorem{theorem}{Theorem}[section]
\newtheorem{lemma}[theorem]{Lemma}
\newtheorem{proposition}[theorem]{Proposition}
\newtheorem{corollary}[theorem]{Corollary}
\theoremstyle{definition}
\newtheorem{definition}[theorem]{Definition}
\theoremstyle{remark}
\newtheorem{remark}[theorem]{Remark}

\newcommand{\code}[1]{\texttt{#1}}
\newcommand{\fl}{\operatorname{fl}}
\newcommand{\df}{\operatorname{df}}

% The measured figures, from docs/experiments/keyword-and-ledger-2026-11-11.json.
\newcommand{\nItems}{400{,}274}
\newcommand{\nUnits}{1{,}369{,}944}
\newcommand{\nPostings}{117{,}975{,}818}
\newcommand{\nRecords}{387{,}838}
\newcommand{\nPaths}{400}
\newcommand{\nBatch}{2{,}050}
\newcommand{\nQueries}{16}
\newcommand{\nAnswers}{65}
\newcommand{\kwBefore}{14.2}
\newcommand{\kwBeforeMax}{33.4}
\newcommand{\kwAfter}{1.20}
\newcommand{\kwAfterMax}{5.87}
\newcommand{\pathsBefore}{5.47}
\newcommand{\pathsAfter}{5.5}
\newcommand{\ledgerBefore}{2.47}
\newcommand{\ledgerAfter}{7}
\newcommand{\singleBefore}{7.7}
\newcommand{\singleAfter}{0.027}
\newcommand{\migrationLex}{115 to 122}

\title{Keyword Search and Ledger Lookups at Corpus Scale:\\
Kept Statistics, One Statement, and an Index by Output Hash}
\author{Alberto Serrano-Calva}
\date{November 11, 2026}

\begin{document}
\maketitle

\begin{abstract}
A knowledge store answers keyword queries from a BM25 index kept in SQLite and links every
attestation of its append-only ledger to the attestations that produced its inputs. On the
store's corpus instance, \nItems{} items in \nUnits{} lexical units with \nPostings{} postings and
\nRecords{} attestations, a title query under the class filter the store always applies took a
warm median of \kwBefore{}~s and up to \kwBeforeMax{}~s, the store's own query shape (a class and
\nPaths{} paths) \pathsBefore{}~s, and the producers of one digest \ledgerBefore{}~s, a lookup an
item landed on its own asks three times. Neither cost is in the data. With no planner statistics,
SQLite counted each query term's document frequency by visiting every unit, and the ledger was
decoded record by record in Python. We keep the statistics instead of counting them: per class,
the current units and their total length, and per class and term the current units holding the
term, moved in the transaction of every write that changes which units are current, so a reader's
statistics remain the sums of its cleared classes' rows and a filtered read remains a function of
the cleared rows alone. Scoring is one statement that starts from the query terms' postings or
from the given paths' units and writes the sum over the terms out in their order, since SQLite's
own sum compensates its rounding: the scores are the old statements' to the last bit. The ledger
gains a table from each output hash to its record and action, written by a trigger in the
append's own statement. On the same clone the warm medians become \kwAfter{}~s, \pathsAfter{}~ms
and \ledgerAfter{}~\textmu s, items landed one at a time go from \singleBefore{}~s to
\singleAfter{}~s an item, and every one of the \nAnswers{} answers compared is identical to the
one before. Property tests check the identity, the statistics against a recount after any
sequence of writes, the clearance rule, and the ledger table against a scan.
\end{abstract}

\section{Introduction}\label{sec:intro}

The store this article concerns keeps two indexes that grow with everything it lands. The lexical
index holds the chunks of every landed item as BM25 postings \citep{robertson2009bm25}, one row a
term a unit, beside each unit's path, version, class, length and a flag that marks the current
version of its path; a lexical \code{near} reads it, and so does the check that tells a landing
whether an item is already held. The attestation ledger records every action with the hashes of
its inputs and outputs; an attestation's parents are the attestations that produced its inputs,
found by asking which records name a hash among their outputs, and a landing asks whether its
digest is attested already before it attests it.

A corpus instance was built to find where the store stops scaling. At its first stage it holds
\nItems{} items, every one of them in the lexical index (\nUnits{} units, \nPostings{} postings)
and in the ledger (\nRecords{} records). Both questions above were past any budget there. A keyword
query took seconds to tens of seconds, and its cost was the number of units, not the number of
postings it read; the producers of one digest took \ledgerBefore{}~s, and a landing that goes item
by item asks three times an item. This article removes both costs without changing an answer.
The contributions:

\begin{enumerate}
\item Statistics kept, not counted (\cref{sec:theory,sec:impl}): per class and per class and term,
  moved by the net currency changes of each transaction, read before any posting of a changed unit
  is deleted, with a proof that the tables equal a recount after every committed transaction and
  that a reader's statistics depend on its cleared classes' units alone.
\item One statement (\cref{sec:impl}): the arithmetic, the sum over the terms and the top $k$ in
  SQL, from the postings or from the paths' units, with the sum written out in the reference's
  order, which makes the result the reference's to the bit where SQLite's \code{sum()} does not.
\item An index of the ledger by output hash, kept by a trigger, so that every writer, of any
  version, keeps it in its own statement.
\item A stress battery that runs on a copy-on-write clone of an instance's stores, and the
  measurements before and after the change (\cref{sec:eval}).
\end{enumerate}

\section{Background and related work}\label{sec:background}

BM25 scores a unit $d$ for a query $q$ as a sum over the query's distinct terms,
\[
  s(q, d) = \sum_{t \in q} \mathrm{idf}(t)\,
            \frac{\mathit{tf}(t, d)\,(k_1 + 1)}{\mathit{tf}(t, d) + k_1 (1 - b + b\,\ell(d)/\bar\ell)},
  \qquad
  \mathrm{idf}(t) = \ln\!\Big(1 + \frac{n - \df(t) + 0.5}{\df(t) + 0.5}\Big),
\]
with $k_1 = 1.2$ and $b = 0.75$, where $n$ is the number of units, $\bar\ell$ their mean length and
$\df(t)$ the units holding $t$ \citep{robertson2009bm25}. Every quantity but $\mathit{tf}$ and
$\ell(d)$ is a statistic of the collection, and in this store the collection is a reader's: the
current units of the classes it is cleared for. An inverted index serves the sum by reading the
postings of the query's terms \citep{zobel2006inverted}, term at a time or document at a time, and
pruning methods skip postings that cannot reach the top $k$ \citep{turtle1995query}.

A relational engine chooses the order of a join from statistics of the tables
\citep{selinger1979access}. The store's SQLite file had none, and without them the engine took an
equality on the index \code{(current, provenance)} as being as selective as an equality on a term:
it counted a term's document frequency by visiting every current unit and probing the postings for
each, a million probes a term whatever the term. Statements can name their order instead
(\code{CROSS JOIN}, which SQLite never reorders) and their index (\code{INDEXED BY}), so the plan
does not depend on what the planner believes.

A sum of floating-point numbers depends on the order of its terms
\citep{ieee754-2019,higham1993accuracy}. Compensated summation carries the rounding error of each
addition into the next \citep{kahan1965further,neumaier1974rundungsfehler}: it is more accurate,
and for that reason it is a different function of the same terms. SQLite's \code{sum()} and
\code{total()} return the compensated sum (\cref{sec:eval} shows a case), so a score summed by
them is not the score the store's Python loop computed.

The store's clearance rule is a noninterference requirement \citep{goguen1982security}: what a
reader sees must be a function of what it is cleared to see. For BM25 it reaches the statistics: a
filter applied after scoring with every class's statistics would let the scores reveal uncleared
units, which is why engines that keep one set of statistics for everything they hold do not meet
it unless given a table a cleared set. Kept statistics meet it when they are kept per class and a
reader sums its own classes' rows. The upkeep must be atomic with the writes it follows
\citep{haerder1983principles}: a transaction that changes the units and dies before the
statistics must leave neither.

\section{The kept statistics and the written-out sum}\label{sec:theory}

\begin{definition}[Statistics]
A store holds units $u$, each with a class $c(u)$, a length $\ell(u)$, a currency bit
$\kappa(u) \in \{0, 1\}$ and postings $P(u)$, a set of terms. For a class $c$ and a term $t$,
\[
  N_c = \textstyle\sum_{u : c(u) = c} \kappa(u), \quad
  L_c = \textstyle\sum_{u : c(u) = c} \kappa(u)\,\ell(u), \quad
  D_{c,t} = \textstyle\sum_{u : c(u) = c,\, t \in P(u)} \kappa(u).
\]
The kept tables hold $(c, N_c, L_c)$ for $N_c > 0$ and $(t, c, D_{c,t})$ for $D_{c,t} > 0$. A
reader cleared for a set of classes $C$ scores with $n = \sum_{c \in C} N_c$,
$\bar\ell = \sum_{c \in C} L_c / n$ and $\df(t) = \sum_{c \in C} D_{c,t}$.
\end{definition}

A transaction changes some units' currency, writes new units with their postings (not current), and
may delete a version's units with their postings. Let $\delta(u) = \kappa'(u) - \kappa(u)$ be the
net change of $u$ over the transaction, which is $0$ for a unit made current and then superseded
within it.

\begin{proposition}[The tables are a recount]\label{prop:recount}
Suppose a transaction adds $\delta(u)$ to $N_{c(u)}$, $\delta(u)\,\ell(u)$ to $L_{c(u)}$ and
$\delta(u)$ to $D_{c(u),t}$ for each $t \in P(u)$, reading $P(u)$ while $u$'s postings are still
those it had while current, drops the rows that reach zero, and commits with its writes or not at
all. Then after every committed transaction the tables equal the statistics counted from the
units.
\end{proposition}

\begin{proof}[Sketch]
Each statistic is a sum over units of $\kappa(u)$ times a weight that depends on $u$ alone ($1$,
$\ell(u)$, or the indicator of $t \in P(u)$). A unit whose currency the transaction leaves alone
contributes the same before and after: new units enter with $\kappa = 0$ and their postings, and
postings are deleted only with their unit, after its change has been applied. A changed unit's
contribution changes by $\delta(u)$ times its weight, which is what is added; reading $P(u)$
before any of its postings is deleted makes the weight the one the unit had while current. By
induction on committed transactions, from a store counted once by the migration (or empty after a
reset), the tables equal the recount. A transaction that dies commits neither its units nor its
deltas.
\end{proof}

\begin{proposition}[A reader's statistics are its classes' own]\label{prop:clearance}
For any two stores whose units of the classes in $C$ are the same, a reader cleared for $C$ is
given the same $n$, $\bar\ell$ and $\df$, and so the same scores for every query.
\end{proposition}

\begin{proof}[Sketch]
$N_c$, $L_c$ and $D_{c,t}$ for $c \in C$ are sums over units of class $c$; the reader's statistics
are sums of these alone, and its rows are read from units of classes in $C$.
\end{proof}

The statement adds a unit's per-term weights in the order of the sorted query terms. Writing that
sum out in SQL needs a value for the terms the unit lacks; it is $0.0$.

\begin{proposition}[The written-out sum is the reference's]\label{prop:sum}
Let $w_1, \dots, w_m$ be positive finite doubles and $x_j = w_j$ for $j$ in a set $J$ and
$x_j = 0$ otherwise. Then
$\fl(\cdots\fl(\fl(x_1 + x_2) + x_3)\cdots + x_m)$ equals the left fold over $J$ in increasing
order starting from $0$, $\fl(\cdots\fl(\fl(0 + w_{j_1}) + w_{j_2})\cdots)$.
\end{proposition}

\begin{proof}[Sketch]
IEEE~754 addition is correctly rounded \citep{ieee754-2019}, and the exact sum of a finite $s \ge 0$
and $0$ is $s$, a double; so $\fl(s + 0) = s$, and a term the unit lacks leaves the partial sum as
it is. The partial sums before $j_1$ are $0$, and $\fl(0 + w_{j_1}) = w_{j_1}$. From there the two
sequences add the same terms to the same partial sums in the same order. Each weight is computed
with the reference's operations in the reference's order on the same operands, so it is the same
double.
\end{proof}

\begin{remark}
Compensated summation is not that fold. For $1 + 2^{-53} + 2^{-53}$ the left fold gives $1$
(each addition rounds to even) and the compensated sum gives $1 + 2^{-52}$, which is what SQLite's
\code{sum()} returns over the same rows in the same order (the control test).
\end{remark}

The ledger's table is simpler. A trigger on the ledger inserts, for an appended record $a$ and
each position $j$ of its output hashes, the row $(h_j(a), \mathrm{id}(a), \mathrm{action}(a), j)$,
in the statement that appends $a$; an append that is ignored (the record exists) or rolled back
writes none. So after every committed statement the table is exactly the set of such rows over the
ledger, and each lookup is a function of that set: the producers of a set of hashes are the ids
of its rows with those hashes, the producers by output group the (action, id) pairs of the rows by
hash, one pair a row, as the records list them, and the producers of one action are those pairs
filtered by action, each id once.

\section{Implementation}\label{sec:impl}

\paragraph{The tables and the migration.} Two tables beside the units in the same file hold
the statistics, \code{class\_stats} with a class's units and tokens and \code{term\_stats} with a
term's units in a class, keyed by the class and by the term and class. A store from before them counts them on its first open, in one transaction: the
classes from the units, then for each class one pass over the postings in their key order with
the units joined by id, so the counts arrive grouped by term and no sort is needed. A store whose
units another tokenizer made is not counted: it neither lands nor scores until it is reset and
re-derived, and its reset counts the empty store. A key in the store's \code{meta} table marks the
statistics counted.

\paragraph{The upkeep.} \Cref{alg:upkeep} is the landing's part; a retire is its last two lines.
The currency moves by two statements a version, each returning the ids it changed, and the net
change of each unit is kept in a map for the whole batch, so a unit made current and superseded
within it costs nothing. A version landed again with other chunks deletes its units: the units of
it that were current leave the statistics first, while their postings exist. The deltas are
applied once a batch, as one grouped read of the changed units' postings and one upsert a
(term, class) pair. Row triggers on the units would keep the same tables with no code in the
store, but a row trigger fires for each unit and makes one upsert for each of its postings,
86 a unit on average in the corpus, where the batch makes one for each distinct (term, class)
pair of all its units together; the batch is the design.

\begin{algorithm}[t]
\caption{The statistics' upkeep inside one landing transaction.}\label{alg:upkeep}
\begin{algorithmic}[1]
\Require versions $(p, v, \mathit{class}, \mathit{chunks})$; a map $\Delta$ from unit to net change
\State begin; hold the writer's key in \code{meta}
\For{each version $(p, v)$}
  \If{the units held for $(p, v)$ are not the chunks}
    \State $\Delta \mathrel{-}= $ the units of $(p, v)$ set not current
    \State apply and remove $\Delta$'s entries for those units \Comment{their postings still exist}
    \State delete their postings and units; write the chunks' units (not current) and postings
  \EndIf
  \State $\Delta \mathrel{-}= $ units of $p$'s other versions set not current \Comment{\code{RETURNING}}
  \State $\Delta \mathrel{+}= $ units of $(p, v)$ set current
\EndFor
\State apply $\Delta$: per sign, group the units' postings by (term, class); upsert; drop zeros
\State release the writer's key; commit
\end{algorithmic}
\end{algorithm}

\paragraph{The guard.} The upkeep is this class's; a process that writes the units without it, an
older version of the store still running when the store was migrated, would leave the statistics
wrong and nothing would show it. So once the statistics are counted, three triggers on the units
refuse an insert, a currency or class change, and a delete unless the writing transaction holds a
key in \code{meta}, which the class's transactions write at their start and remove before they
commit, so no other connection ever sees it. The postings need no guard: they are written and
deleted only with their unit. The cost is one lookup of a one-row table a changed unit.

\paragraph{The statement.} A query's terms with a document frequency above zero go into a
temporary table with their positions $i$ in sorted order and their weights $\mathrm{idf}$; the
paths, when given, into another. The rows are read postings first,
\code{q CROSS JOIN postings p ON p.term = q.term CROSS JOIN units u ON u.unit\_id = p.unit\_id},
or paths first, \code{paths w CROSS JOIN units u INDEXED BY units\_path ON u.source\_path = w.path
AND u.current = 1}, then the postings by (term, unit). Each row's weight is
\code{q.w * ((p.tf * ?) / (p.tf + ? * (? + (? * u.length) / ?)))}, with $k_1 + 1$, $k_1$, $1 - b$,
$b$ and $\bar\ell$ bound as doubles computed as the reference computes them; the rows are grouped
by unit, the score is \code{COALESCE(max(c) FILTER (WHERE i = 0), 0.0) + COALESCE(max(c) FILTER
(WHERE i = 1), 0.0) + \dots} in that order (\cref{prop:sum}; a filtered aggregate is not called for
the rows of the other terms, which made the long titles 6 to 27\% faster than the same sum written
with \code{CASE}), and \code{ORDER BY score DESC, unit\_id LIMIT k} returns the top, ties by unit id
as before. \code{scores()}, each path's best unit, is the same statement grouped by path with
\code{max}. SQLite parses the written-out sum as a tree as deep as its terms and refuses a depth
over 1{,}000, so a query of more than 500 distinct terms has the same weights read out of the
statement in the terms' order and summed in Python, the reference's sum again. Without
\code{INDEXED BY}, SQLite drove the paths shape from the class's index and visited every current
unit of the class for each path.

\paragraph{The ledger.} \code{output\_hashes(hash, att\_id, action, pos)}, keyed by all four, is
created, filled from the records by one \code{json\_each} pass and given its trigger in one
transaction on the first open of a ledger without it. The lookups read it by seeks, a statement
for a few hundred hashes, and only read: the statement it replaced filled a temporary table, which
in the ledger's transaction mode opened a transaction that kept a shared lock on the file until the
connection next committed. \code{attested\_by} and \code{already\_attested} read the producers by
output and keep the action asked for; a store without the table, as tests build, is still asked
the old way.

\section{Evaluation}\label{sec:eval}

\paragraph{Method.} \code{eval/stress\_battery.py} clones an instance's stores with \code{cp -c}
(a copy-on-write clone: nothing of the source is opened for writing) and runs a fixed battery:
keyword queries free (the class filter and none) and in the store's call shape (a class and
\nPaths{} random item paths, \code{scores} and \code{search}), warm in one process (the first pass
apart, then three) and in a fresh process a query (three repetitions); the producers of single
digests and of a batch of \nBatch{}; 2{,}000 synthetic items landed through the writer, batched and
lazy, three batches a process in three processes, and three items one at a time; and the drain's
pending set, which no line of the change reads, as a control of the machine. A synthetic item is a
real item of the corpus with one first line added, so its digest is new and its text, chunks and
terms are a real item's. The queries are \nQueries{} titles, ten from an earlier study of the
stores and six drawn independently, with one to nineteen distinct terms and 32 to 6.5 million
postings, named by their sizes and never by their text. The corpus's lexical file predates the
store's record of its tokenizer, so it is opened with the tokenizer that made it. The machine is an
Apple M2 Pro with 12 cores and 32~GB, shared with other work: every timing carries the one-minute
load average beside it, BLAS is capped at four threads, and every process ends itself above
3.8~GB resident. One run before the change, on its parent commit, one after, and the keyword and
ledger parts again on the final commit. The whole-stratum load of a sentence admitting the
dialogue stratum was not run: on the whole corpus it needs tens of gigabytes.

\begin{table}[t]
\centering
\caption{Before and after on clones of the corpus instance: warm medians over the queries of each
query's median, with the range over the queries, and the one-minute load average.}\label{tab:results}
\small
\begin{tabular}{@{}lrrr@{}}
\toprule
measure & before & after & load \\
\midrule
title query, class filter (16) & 14.2 s (1.47--33.4 s) & 1.20 s (0.1 ms--5.87 s) & 2.8--7.4 / 2.2--11.3 \\
title query, no filter (16) & 9.42 s (0.98--23.2 s) & 1.16 s (0.1 ms--5.70 s) & 2.0--3.2 / 2.5--5.6 \\
the store's shape, \code{scores} (5) & 5.47 s (1.22--18.9 s) & 5.5 ms (1.9--12.1 ms) & 1.9--4.5 / 2.5 \\
the store's shape, \code{search} (5) & 8.43 s (2.33--15.7 s) & 4.6 ms (1.8--11.1 ms) & 3.1--4.2 / 2.5 \\
fresh process, 14 terms, class filter & 24.1--24.4 s, 468 MB & 3.95--3.97 s, 74 MB & \\
producers of one digest & 2.47 s & 7 \textmu s & 2.1--3.2 / 2.5--3.7 \\
\code{attested\_by} through the landing & 2.48 s & 8 \textmu s & \\
producers of 2{,}050 digests & 304 ms & 5.9 ms & \\
landing 2{,}000 items, batches 1--2 & 10.40 s & 10.64 s & 2.2--2.8 / 3.4--6.3 \\
three items one at a time & 23.1 s & 0.082 s & \\
the drain's pending set (control) & 2.24 s & 2.33 s & 2.3--2.8 / 3.3--3.4 \\
\bottomrule
\end{tabular}
\end{table}

\paragraph{Results.} \Cref{tab:results}. The title queries under the class filter fall from a
median of \kwBefore{}~s to \kwAfter{}~s and the slowest from \kwBeforeMax{}~s to \kwAfterMax{}~s;
the class-filtered pass of the final run met a load of 6.7 to 11.3 from another job for its
second half, and the unfiltered pass, which scores the same units here, ran at 2.5 to 5.6 and
gives 1.16~s. A free query still reads every posting of its terms, about 0.9~\textmu s a posting,
but no term costs a visit of every unit any more. The store's own call shape falls from seconds
to milliseconds, as does a fresh process's query, whose peak memory falls from 468~MB to 74~MB
with the dictionary of every scored unit the old loop built. The producer lookup falls from
2.47~s to 7~\textmu s and the batch's from 304~ms to 5.9~ms. A batched landing is unchanged within
its spread: its lexical stage's median rose 0.55~s at a load one to four higher, while the ledger's
lookup saved 0.3 to 0.7~s a batch and the trigger cost 27~\textmu s an append. Landing items one at
a time, \singleBefore{}~s an item before, 99\% of it the three passes over the ledger, takes
\singleAfter{}~s. The pending set, which no part of the change reads, took 2.24~s and 2.33~s, so the
runs met the same machine within 4\%. The migrations run once a store: counting the statistics of
the \nUnits{} units took \migrationLex{}~s on cold clones and added 29.6~MB to the 4.6~GB file; the
ledger's table took 1.8 to 2.7~s and added 64~MB, 166~B a record. Every one of the \nAnswers{}
answers the battery keeps (top tens with their scores, path scores, producer sets and lists) is
identical before and after, scores compared as the bits of their doubles, and so are the fresh
processes' answers.

\paragraph{Tests.} Over generated stores of four classes, three tokenizers, shared vocabularies
(so scores tie), empty and absent classes and paths given and not, the statement's top $k$ and
scores equal the old statements' (kept verbatim as the reference) for $k \in
\{-1, 0, 1, 2, 3, 10\}$, and so do they with the sum taken outside the statement (its threshold
lowered to one term) and for a query of 1{,}200 terms; after any sequence of landings, landings
again with other chunks, retires and re-derivations, the tables equal a recount, also after
reopening; and a store that also holds other classes' units, landed and retired around the cleared
ones, gives a cleared reader the statistics, scores and ranking of a store that holds the cleared
units alone. Each runs 60 examples in the suite and passed 1{,}000 to 1{,}500 in a longer run. A store written by the old code counts
its statistics on its first open; a stale one waits for its reset; a write to the units from
another connection is refused, also the old landing; a landing or a retire that dies between the
unit rows and the statistics, by an exception or by the process ending, leaves neither. For the
ledger, after any sequence of appends, appends again, and deferred blocks that commit or die, the
table equals a scan of the records, and the four lookups answer what the scan and the replaced
statement answer, also on a ledger the old code wrote and after other writers append to it.

\section{Limitations}\label{sec:limits}

A free query still reads every posting of its terms, and its time grows with them (\kwAfterMax{}~s
for the query with 6.5 million postings); pruning would read a fraction \citep{turtle1995query}.
The kept statistics are the present's: a read at a past cut must count its own under the cut's mask
or keep them per epoch. The guard turns an older process's writes into errors, by design; the
statistics of a stale store wait for its reset, and the corpus's file, made by the earlier
tokenizer, must be re-derived before the store's current tokenizer scores it. The migration holds
the file's write lock for its \migrationLex{}~s on the corpus, and another writer waiting longer than
the store's 30~s timeout fails. The paths shape reads every unit of every path given: a call with
hundreds of thousands of paths, which a mostly unembedded corpus would make, needs the pending set
redesigned first. The corpus holds one class, so its readers' sums are over one row; the property
tests hold four. The measurements are one run before and one after on one shared machine; the
cold behaviour of a file larger than memory was not measured.

\section{Future work}\label{sec:future}

The posting lists as arrays in immutable segments, with the statistics kept per class and per
segment and an exact pruned search, are the change that makes a free query's time follow the
postings it needs rather than the postings its terms hold. The pending set as a mask over one
ranking removes the paths shape's growth. The upkeep's invariant (\cref{prop:recount}) and the
clearance rule (\cref{prop:clearance}) have the shape of the store's other laws and can be stated
over a finite model and proved, beside the property tests that check them now.

\section{Conclusion}\label{sec:conclusion}

The two slowest reads of the corpus instance were not slow because of the data but because of
what was recomputed for each question: a document frequency counted by visiting every unit, a
ledger decoded record by record. Keeping the statistics with the writes that change them, naming
the statements' order, writing the sum out in the reference's order and indexing the ledger by
output hash make the same answers cost milliseconds, with the clearance rule and the answers as
they were.

\bibliographystyle{plainnat}
\bibliography{refs}
\end{document}
